\documentclass[conference]{IEEEtran}

\usepackage{amsmath,amssymb}
\usepackage{graphicx}
\usepackage{booktabs}
\usepackage{array}
\usepackage{cite}
\usepackage{url}
\usepackage{algorithm}
\usepackage[noend]{algpseudocode}
\usepackage[hidelinks]{hyperref}

\newcommand{\meanpm}[2]{#1\,$\pm$\,#2}
\newcommand{\protocol}[1]{\textsf{#1}}
\newcommand{\cytofair}{\textsc{CytoFair}}
\newcommand{\placeholder}[1]{\texttt{[#1]}}
\setlength{\textfloatsep}{5pt plus 1pt minus 2pt}
\setlength{\dbltextfloatsep}{5pt plus 1pt minus 2pt}
\setlength{\floatsep}{4pt plus 1pt minus 2pt}
\setlength{\dblfloatsep}{4pt plus 1pt minus 2pt}
\setlength{\intextsep}{4pt plus 1pt minus 2pt}
\setlength{\abovecaptionskip}{2.5pt}
\setlength{\belowcaptionskip}{0pt}
\renewcommand{\arraystretch}{0.9}

\title{Do We Train on Test Cells? Provenance Leakage in Public Cervical Cytology Benchmarks}

\author{
\IEEEauthorblockN{Snehasish Ghosh\IEEEauthorrefmark{1},
Soham Mitra\IEEEauthorrefmark{2}, and
Nibaran Das\IEEEauthorrefmark{2}}
\IEEEauthorblockA{\IEEEauthorrefmark{1}Indian Institute of Science, Bengaluru, India\\
Email: snehasishg@iisc.ac.in}
\IEEEauthorblockA{\IEEEauthorrefmark{2}Department of Computer Science and Engineering,
Jadavpur University, Kolkata, India\\
Emails: \placeholder{SOHAM EMAIL}, \placeholder{NIBARAN EMAIL}}
}

\begin{document}
\maketitle

\begin{abstract}
Public cervical cytology benchmarks are built by cropping many cells from far fewer microscopic fields, and the literature splits them at the crop level.
We recover the concealed hierarchy and measure what that costs.
Under the standard stratified 80:20 crop split of SIPaKMeD, \meanpm{92.62}{0.40}\% of test crops have a training parent and \meanpm{52.62}{1.97}\% share literal source pixels with a training crop; Cx22 and CRIC show the same pattern, while single-cell Herlev returns exactly \meanpm{0.00}{0.00}\% and so serves as a negative control on the measurement itself.
The shortcut is learnable: a parent-identity probe reaches \meanpm{25.73}{0.66}\% on a 424-way task against a 0.2358\% chance baseline and a \meanpm{0.17}{0.24}\% label-shuffled comparator, and still reaches \meanpm{18.68}{1.25}\% from inpainted background alone---so the fingerprint is acquisition context, not cell morphology.
Correcting the evaluation unit lowers accuracy for all six architectures tested (mean 1.37 points), exceeding the 0.62-point spread between the two strongest published five-class results, though our five-seed design is underpowered to certify individual contrasts after correction.
We release \cytofair: exact provenance maps, fixed parent-disjoint splits, and the audit toolkit, at \placeholder{REPOSITORY URL}.
\end{abstract}

\begin{IEEEkeywords}
cervical cytology, data leakage, provenance, grouped splitting, benchmark audit, reproducibility
\end{IEEEkeywords}

\section{Introduction}
Cervical cancer caused an estimated 660,000 new cases and 350,000 deaths worldwide in 2022~\cite{who2025hpv}, and computer-assisted cytology screening is attractive as a way to support scarce expert attention.
That promise rests on an elementary validity condition: the held-out sample should represent genuinely unseen biological or acquisition units.
The image presented to a network is not always that unit.
Public cervical benchmarks are overwhelmingly \emph{cropped-from-parent} constructions---many cell images derived from far fewer microscopic fields.
Sibling crops from one field inherit stain batch, illumination, focus, background texture, preparation artifacts, and sometimes overlapping pixels.
A crop-level split then places siblings on both sides of evaluation while appearing formally disjoint at the filename level, because parent identifiers are not always carried into the released crop folders.

Whether this matters depends on its scale relative to the differences the field publishes about, and that comparison is unusually stark here.
DeepPap reported 98.3\% on Herlev~\cite{zhang2017deeppap}; DeepCervix reported 99.14\% five-class accuracy on SIPaKMeD~\cite{rahaman2021deepcervix}; a knowledge-distillation framework reported 98.52\% on the same task~\cite{gao2022knowledge}.
The two strongest five-class SIPaKMeD results we cite therefore differ by \textbf{0.62 percentage points}.
Any protocol artifact of comparable magnitude is not a footnote; it is larger than the quantity being contested.
These methods use standard backbones~\cite{he2016resnet,huang2017densenet,sandler2018mobilenetv2,tan2019efficientnet,dosovitskiy2021vit} on releases that made comparison possible: SIPaKMeD provides 4,049 crops from 966 cluster fields~\cite{plissiti2018sipakmed}, Cx22 provides 1,320 generated regions~\cite{liu2022cx22}, CRIC provides 11,534 cells in 400 images~\cite{rezende2021cric}, and Herlev provides 917 single-cell images~\cite{jantzen2005pap}.
Performance engineering has advanced faster than auditing of the sampling hierarchy beneath them.

We recover that hierarchy exactly where the release permits, measure what crop splitting shares, test whether the shared signal is learnable, and build partitions that keep a parent wholly in one arm.
We separate two claims throughout, because only one is established here: recoverable parent information shows a shortcut is \emph{available}, not that every model exploits it.
Our contributions are fivefold:
\begin{itemize}\itemsep0pt
    \item \cytofair, a framework for exact crop-to-parent recovery, fail-closed validation, and provenance-aware split generation;
    \item the first complete dual-branch provenance map for all 4,049 SIPaKMeD crops---two independent recoveries agreeing on 4,049/4,049---plus source recovery for Cx22 and structural audits of CRIC and Herlev;
    \item a structural negative control: single-cell Herlev returns exactly zero on every leakage endpoint, validating the instrument and not only the datasets under test;
    \item measurement of parent straddling, literal cross-split pixel sharing, and a five-seed parent-identity probe with four transformation controls and a shuffled-label comparator, localizing the fingerprint to acquisition context;
    \item fixed parent-disjoint splits, released, with a deliberately bounded report of accuracy sensitivity on SIPaKMeD---the one dataset given the full six-architecture rescore.
\end{itemize}

\section{Related Work}
Our design is assembled from four established lines rather than proposed from scratch.
Barz and Denzler supply the template: audit a benchmark for train--test near-duplicates, quantify the exposure, and release a corrected evaluation set rather than only a complaint, as they did for CIFAR-10/100 (3.3\% and 10\% of test images) with ciFAIR~\cite{barz2020cifair}; we substitute exactly recovered provenance for image similarity as the grouping criterion, reserving similarity for a labeled sensitivity analysis.
Laroca \emph{et al.} supply the re-evaluation protocol, re-running six recognition models under duplicate-free license-plate splits~\cite{laroca2023lpr}; our multi-backbone protocol comparison is the same manoeuvre in a medical setting.
Packh\"auser \emph{et al.} supply the mechanism instrument: their demonstration that chest radiographs support patient re-identification at 95.55\% accuracy~\cite{packhauser2022reidentification} establishes identity probing as evidence that a nuisance signal is present and learnable independent of any downstream task, which Sec.~\ref{sec:probe} adapts from patient to source-field identity and extends with a transformation ablation.
Kapoor and Narayanan supply the taxonomy situating train--test non-independence within a broader account of leakage in ML-based science~\cite{kapoor2023leakage}.
Crop-level cervical accuracies are not imported into our comparison, because different splits do not define comparable tests; we use only their \emph{spread}.
The remaining gap is an exact, validated provenance audit and a corrected public evaluation resource for cervical-cytology crop benchmarks.

\section{Method}
\cytofair{} treats provenance recovery as a prerequisite, not a post-hoc clustering step: released or exactly reconstructed parent identity defines the independent grouping unit, while image-derived similarity affects only the conservative \protocol{P3} sensitivity split.
Algorithm~\ref{alg:cytofair} states the full procedure.

\subsection{Exact Provenance Recovery}
For each released crop $x_i$, let $g(x_i)=p_j$ denote its source parent.
SIPaKMeD Branch A parses the released filename hierarchy; Branch B ignores that assignment entirely, searches all same-class parent fields, and selects the contained region maximizing normalized cross-correlation.
Branch B is not a convenience---it is why the map cannot be dismissed as an artifact of filename parsing.
Recovery is accepted only if the branches agree on at least 99\% of common crops, every bounding box lies inside its parent, every crop has exactly one parent, crop totals equal the release, and reconstructed class-parent counts equal the published vector.
Cx22 parents and requested regions are decoded from the authors' MATLAB generator metadata and checked against the 686-image CCEDD source release; CRIC parent identities come from its released classifications table; one Herlev crop is one parent.

\subsection{Structural and Literal Sharing}
For a parent with $k$ crops assigned independently to 80\% training and 20\% test, the probability of straddling both arms is
\begin{equation}
P_{\mathrm{straddle}}(k)=1-0.8^k-0.2^k,
\label{eq:straddle}
\end{equation}
and the analytic dataset expectation is its mean over parents---a prediction the measurement must reproduce.
For an allocation with training-parent set $G_{\mathrm{tr}}$ and test crops $T$,
\begin{equation}
L_{p}=\frac{|G_{\mathrm{tr}}\cap G_{\mathrm{te}}|}{|G|},\quad
L_{c}=\frac{1}{|T|}\sum_{x\in T}\mathbb{1}[g(x)\in G_{\mathrm{tr}}].
\label{eq:leakage}
\end{equation}
For SIPaKMeD the recovered bounding boxes additionally permit literal sharing: $L_{\mathrm{pix}}$ is the fraction of test crops whose parent has at least one training crop with nonempty source-coordinate intersection.
This endpoint requires no model and admits no interpretive latitude---crops either share pixels or they do not.
Observed five-seed estimates are checked against an exact 10,000-allocation Monte Carlo interval under the implemented fixed-size, class-stratified rule.

\subsection{Parent-Identity Mechanism Probe}
\label{sec:probe}
Parents with at least four crops enter a balanced closed-set task, each contributing equal training and test crops, and a ResNet-18 is trained with cross-entropy over parent labels.
The question is whether two \emph{different} crops from one field retain a learnable common signature after the cellular content has changed; an affirmative answer establishes availability of a shortcut, not causal use by a class classifier.
Four prespecified transformations isolate candidate substrates: grayscale removes hue, channel-statistics matching follows Reinhard \emph{et al.}~\cite{reinhard2001color}, background-only masks use Otsu thresholds~\cite{otsu1979threshold} with Telea inpainting~\cite{telea2004inpainting}, and high-pass filtering suppresses low-frequency illumination.
The background-only condition is decisive: it removes the cell.
Images, optimizer, epochs, and seeds are retained for a parent-label-shuffled control, whose accuracy must fall to chance if the probe measures what it claims.

\subsection{Matched Protocols and Statistics}
\protocol{P1} is an 80:20 class-stratified crop split---the protocol the literature uses.
\protocol{P2} assigns complete recovered parents to one arm.
\protocol{P3} joins parents connected by cross-parent candidates satisfying both 64-bit pHash and dHash Hamming thresholds of six, then groups the entire connected component; it is a sensitivity analysis, not a claim that every candidate is a proven duplicate.
\protocol{P1$'$} retains the \protocol{P1} test set and downsamples only its training arm to the exact \protocol{P2} training count, so accuracy change cannot be attributed to training-set size; on SIPaKMeD the counts are already identical, so \protocol{P1$'$}=\protocol{P1} and it is omitted from tables.
Classifiers are ImageNet-pretrained ResNet-18/50~\cite{he2016resnet}, DenseNet-121~\cite{huang2017densenet}, MobileNet-V2~\cite{sandler2018mobilenetv2}, EfficientNet-B0~\cite{tan2019efficientnet}, and ViT-B/16~\cite{dosovitskiy2021vit}, with only final heads replaced and identical seeds paired across protocols.
Normal paired differences receive a two-sided paired $t$-test with Cohen's $d_z$; the preregistered fallback is Wilcoxon with rank-biserial correlation.
All \protocol{P1}$-$\protocol{P2} and \protocol{P1$'$}$-$\protocol{P2} tests form one 12-test Benjamini--Hochberg family~\cite{benjamini1995fdr}, with 10,000 paired bootstrap resamples for effect-size intervals.
Because \protocol{P1$'$} duplicates \protocol{P1} here, that family contains exact duplicates; duplicating every hypothesis scales both family size and rank proportionally, so adjusted values equal those of the six-test family.

\begin{algorithm}[t]
\caption{Fail-closed provenance-aware evaluation}
\label{alg:cytofair}
\small
\begin{algorithmic}[1]
\Require Released observations $X$, source fields $S$, seeds $\mathcal{R}$
\State Recover $g_A:X\!\to\!S$ from released identifiers
\State Recover $g_B:X\!\to\!S$ by independent content matching
\If{agreement, containment, uniqueness, or count gates fail}
    \State \textbf{stop} and emit an audit ledger
\EndIf
\ForAll{$r\in\mathcal{R}$}
    \State Build \protocol{P1}, \protocol{P1$'$}, \protocol{P2}, \protocol{P3}; measure $L_p$, $L_c$, $L_{\mathrm{pix}}$
    \State Train the fixed probe and classifier matrix
\EndFor
\State Report mean $\pm$ SD, paired effects, adjusted inference
\State Export provenance, split hashes, and raw runs
\end{algorithmic}
\end{algorithm}

\section{Experiments}
\subsection{Datasets, Integrity, and Setup}
Branch A and Branch B agree on 4,049/4,049 SIPaKMeD crops, all recovered boxes are contained, and the class-parent vector is 126/108/238/271/223, matching the published release; no observation is excluded.
The 4,049 crops map to 966 parents, Cx22's 1,320 regions to 514 recovered CCEDD sources, CRIC's 11,534 cells to 400 images, and Herlev contains 917 one-crop parents.
The crop-to-parent ratio drives everything downstream and is uneven: SIPaKMeD crop counts are nearly balanced across classes (787--831) while parent counts range from 108 Parabasal to 271 Metaplastic fields, on a 3,239/810 train/test allocation.
All models use AdamW, learning rate and weight decay $10^{-4}$, batch size 64, $224\times224$ RGB inputs, resized-crop and H/V-flip augmentation, and ImageNet initialization; classifiers train 20 epochs and the probe 15.
Every stochastic cell uses seeds 17, 29, 43, 71, and 101, reported as mean $\pm$ sample SD, on at most three NVIDIA RTX A6000 GPUs (one per run) with PyTorch 2.9.0 and CUDA 12.8.
Accuracy is the preregistered primary endpoint; macro-F1 and balanced accuracy are retained concordance metrics.

\subsection{Result 1: The Crop Is Not the Independent Unit}
Table~\ref{tab:structural} is the principal result, and it is close to deterministic.
Observed straddle means fall inside their exact Monte Carlo intervals for all four datasets and reproduce the analytic prediction of~\eqref{eq:straddle}: SIPaKMeD predicts 44.83\% and measures \meanpm{45.63}{1.04}\%, Cx22 predicts 33.93\% and measures \meanpm{34.01}{1.10}\%, CRIC predicts 94.88\% and measures \meanpm{94.35}{1.10}\%.
Mechanism and measurement agree to within sampling error on all three.

The crop-level consequences are severe.
Under \protocol{P1}, \meanpm{92.62}{0.40}\% of SIPaKMeD test crops have a parent in training; Cx22 \meanpm{85.23}{3.51}\%; CRIC \meanpm{99.97}{0.08}\%.
Among 13,064 within-parent SIPaKMeD crop pairs, 2,185 have intersecting source coordinates, and under \protocol{P1} \meanpm{52.62}{1.97}\% of test crops share literal source pixels with a training crop.
More than half the nominal test set is, in part, training data.

Herlev anchors the interpretation: the only audited benchmark with one cell per field returns exactly \meanpm{0.00}{0.00}\% on every endpoint.
The instrument does not manufacture leakage where the construction cannot produce it, the highest-crop-per-field construction (CRIC) produces the highest straddle rate, and dual-branch recovery plus the analytic check are mutually independent safeguards against a split-implementation artifact.
Both grouped protocols eliminate all measured parent and pixel sharing.

\begin{table*}[t]
\caption{Structural leakage under crop-level and grouped protocols. Stochastic entries are mean $\pm$ sample SD over five fixed seeds. Herlev is the structural negative control.}
\label{tab:structural}
\centering
\scriptsize
\setlength{\tabcolsep}{5.2pt}
\begin{tabular}{lrrrrrrr}
\toprule
Dataset & Crops & Parents & Analytic \protocol{P1} straddle & Observed \protocol{P1} straddle & \protocol{P1} test crop w/ train parent & \protocol{P2} straddle & \protocol{P3} straddle \\
\midrule
SIPaKMeD & 4,049 & 966 & 44.83\% & \meanpm{45.63}{1.04}\% & \meanpm{92.62}{0.40}\% & \meanpm{0.00}{0.00}\% & \meanpm{0.00}{0.00}\% \\
Cx22 & 1,320 & 514 & 33.93\% & \meanpm{34.01}{1.10}\% & \meanpm{85.23}{3.51}\% & \meanpm{0.00}{0.00}\% & \meanpm{0.00}{0.00}\% \\
CRIC & 11,534 & 400 & 94.88\% & \meanpm{94.35}{1.10}\% & \meanpm{99.97}{0.08}\% & \meanpm{0.00}{0.00}\% & \meanpm{0.00}{0.00}\% \\
Herlev & 917 & 917 & 0.00\% & \meanpm{0.00}{0.00}\% & \meanpm{0.00}{0.00}\% & \meanpm{0.00}{0.00}\% & \meanpm{0.00}{0.00}\% \\
\bottomrule
\end{tabular}
\end{table*}

\subsection{Result 2: The Shortcut Is Learnable, and It Is Not the Cell}
On the balanced 424-way task (Table~\ref{tab:probe}), parent identity is recovered at \meanpm{25.73}{0.66}\%---109.1$\times$ the 0.2358\% chance baseline.
The shuffled-label control collapses to \meanpm{0.17}{0.24}\%, with a five-seed 95\% interval of $-0.14$\% to $0.47$\% covering the exact chance value, confirming that the probe measures field identity rather than an artifact of task construction.
Parent information survives every prespecified transformation, and the decisive condition is background-only at \meanpm{18.68}{1.25}\%---79.2$\times$ chance from an image whose cell has been masked out and inpainted.

The fingerprint is therefore not cellular morphology but \emph{acquisition context}: slide preparation, illumination, focus, and debris field.
That is the generalizing observation of this paper.
Because the substrate is context rather than content, the exposure is a property of the \emph{construction}---any crop-level split of any cropped-from-field cytology release carries it, however carefully the class labels were assigned.
Persistence across all four transformations shows the signal is distributed rather than confined to one channel, and does not identify a single exclusive substrate.

\begin{table}[t]
\caption{Parent-ID probe ablation, balanced 424-way task. Chance is 0.2358\%.}
\label{tab:probe}
\centering
\scriptsize
\setlength{\tabcolsep}{4.2pt}
\begin{tabular}{lrr}
\toprule
Input condition & Accuracy & Multiple of chance \\
\midrule
RGB & \meanpm{25.73}{0.66}\% & 109.1$\times$ \\
Grayscale & \meanpm{13.80}{1.12}\% & 58.5$\times$ \\
Stain-normalized & \meanpm{17.31}{0.64}\% & 73.4$\times$ \\
Background only & \meanpm{18.68}{1.25}\% & 79.2$\times$ \\
High-pass & \meanpm{16.27}{0.60}\% & 69.0$\times$ \\
Shuffled labels & \meanpm{0.17}{0.24}\% & 0.7$\times$ \\
\bottomrule
\end{tabular}
\end{table}

\subsection{Result 3: Bounded Accuracy Sensitivity}
Table~\ref{tab:joint} rescores six architectures on SIPaKMeD under identical training settings and identical 3,239/810 split sizes.
Correcting the evaluation unit lowers accuracy for \emph{every} architecture, from \meanpm{0.72}{0.83} points (ResNet-50) to \meanpm{2.25}{1.63} (MobileNet-V2), mean 1.37.
Because \protocol{P1$'$}=\protocol{P1} here, these differences cannot be attributed to reduced training-set size; train and test counts are held exactly fixed across protocols.
Effect sizes are moderate to large ($d_z$ 0.588--1.561) and five of six bootstrap intervals for $d_z$ exclude zero, but with five seeds and a declared 12-test family no contrast survives adjustment; the smallest adjusted value is $q=0.1109$.
The design is therefore \emph{underpowered to certify individual architecture-wise contrasts}---a property of sample size and correction breadth, not evidence of absence.
We report direction and scale, and do not claim a confirmatory universal inflation effect.

\begin{table*}[t]
\caption{SIPaKMeD test accuracy and preregistered paired inference. All models were rerun here under identical 3,239/810 counts; mean $\pm$ sample SD over five seeds. $\Delta$ is \protocol{P1}$-$\protocol{P2} in percentage points; $q$ is BH-adjusted over the declared 12-test family.}
\label{tab:joint}
\centering
\scriptsize
\setlength{\tabcolsep}{5.5pt}
\begin{tabular}{lrrrrrrr}
\toprule
Architecture & \protocol{P1} crop & \textbf{\cytofair{} \protocol{P2}} & \protocol{P3} component & $\Delta$ & Raw $p$ & BH $q$ & $d_z$ \\
\midrule
DenseNet-121 & \meanpm{98.12}{0.28}\% & \textbf{\meanpm{96.72}{0.69}\%} & \meanpm{97.38}{0.50}\% & \meanpm{1.41}{0.90} & 0.0251 & 0.1109 & 1.561 \\
EfficientNet-B0 & \meanpm{97.78}{0.49}\% & \textbf{\meanpm{96.69}{1.49}\%} & \meanpm{97.41}{0.71}\% & \meanpm{1.09}{1.85} & 0.2591 & 0.2591 & 0.588 \\
MobileNet-V2 & \meanpm{97.41}{0.45}\% & \textbf{\meanpm{95.16}{1.50}\%} & \meanpm{96.20}{0.93}\% & \meanpm{2.25}{1.63} & 0.0370 & 0.1109 & 1.377 \\
ResNet-18 & \meanpm{97.19}{0.59}\% & \textbf{\meanpm{96.00}{0.89}\%} & \meanpm{96.35}{1.28}\% & \meanpm{1.19}{1.28} & 0.1081 & 0.1511 & 0.922 \\
ResNet-50 & \meanpm{97.51}{0.47}\% & \textbf{\meanpm{96.79}{0.42}\%} & \meanpm{96.27}{0.53}\% & \meanpm{0.72}{0.83} & 0.1260 & 0.1511 & 0.863 \\
ViT-B/16 & \meanpm{97.14}{0.95}\% & \textbf{\meanpm{95.56}{0.84}\%} & \meanpm{95.53}{1.68}\% & \meanpm{1.58}{1.70} & 0.1059 & 0.1511 & 0.931 \\
\bottomrule
\end{tabular}
\end{table*}

\subsection{Registered Predictions and What They Rule Out}
Two preregistered predictions were not supported, and both are retained.
First, the monotonic ordering \protocol{P3}$\leq$\protocol{P2}$\leq$\protocol{P1} fails for four of six architectures.
On inspection the prediction was mis-specified rather than falsified: \protocol{P3} merges parents into similarity-connected components and therefore changes grouped \emph{test composition}, so it is not nested within \protocol{P2} and no monotone finite-sample ordering should have been expected.
Read as measurement rather than failed test, it carries usable guidance---similarity purging on top of parent grouping does not further reduce accuracy, indicating the pHash/dHash components largely capture within-parent structure that \protocol{P2} has already absorbed.
For practitioners building corrected splits, grouping by recovered parent appears sufficient and the more expensive purge is optional.
Second, ranking stability cannot be resolved at this sample size: Kendall's $\tau$ between \protocol{P1} and \protocol{P2} orderings is 0.467 with 4/15 pairwise reversals, but the seed-bootstrap 95\% CI spans 0.071--0.733 with unadjusted $p=0.2722$, indistinguishable from chance across six architectures, and we draw no conclusion from it.
Finally, the frozen run schema retained aggregate metrics and epoch histories but not per-sample predictions or checkpoints, so a post-hoc class confusion matrix would require new inference after the declared hard stop; we neither fabricate one nor infer class-pair errors from aggregate metrics.

\section{Discussion}
The finding is structural rather than architectural.
For three of four audited benchmarks the released image is not the independent sampling unit, and this is now a recovered, doubly validated fact: two independent recoveries agree on 4,049/4,049 crops, observed straddle rates reproduce an analytic prediction, and a single-cell control returns exactly zero.
Under the protocol the literature uses, most SIPaKMeD test crops have a training parent and more than half share literal source pixels with one.
The linking signal is learnable at 109$\times$ chance, and at 79$\times$ chance after the cell itself is removed.

The natural objection is that a 1.37-point mean accuracy change is small.
It is small absolutely and large relative to what the field is contesting.
The two strongest published five-class SIPaKMeD results differ by 0.62 points~\cite{rahaman2021deepcervix,gao2022knowledge}; our measured protocol effect exceeds that spread for \emph{every} architecture tested---the smallest, ResNet-50 at \meanpm{0.72}{0.83}, still exceeds it---and reaches 2.25 points, roughly 3.6$\times$ the published spread, for MobileNet-V2.
This is a comparison of \emph{magnitudes}, not of accuracies: we deliberately do not place our numbers beside published numbers, because different splits do not define comparable tests.
What is comparable is the sensitivity of a fixed pipeline to split protocol, set against the size of the differences the literature reports.
On that comparison, protocol choice is of the same order as, or larger than, the methodological progress being claimed.

We show a shortcut is available, not that classifiers use it.
The probe demonstrates that field identity is recoverable from a crop; it does not demonstrate that a class model routes its prediction through that identity, and our frozen evidence cannot settle it.
Establishing consequence, as distinct from availability, requires a causal design such as intervention on the fingerprint substrate.
One design lesson is recorded rather than repaired: the scientific question was singular, but the preregistered analysis tested it six times and then paid a six-way multiplicity penalty for asking once per backbone.
A pooled omnibus endpoint treating architecture as a nuisance factor would have addressed the same question at materially higher power; we did not preregister one, and do not substitute one after seeing results.

\section{Limitations}
Only SIPaKMeD received the six-architecture rescore; Cx22, CRIC, and Herlev receive structural audits alone, so magnitude claims are restricted to one benchmark.
The per-architecture inference is underpowered at five seeds under a 12-test family, and the accuracy arm therefore estimates rather than certifies.
Interpretation of the weakest cells, MobileNet-V2 under \protocol{P2} and ViT-B/16 under \protocol{P3}, must remain aggregate given the frozen artifacts.
No prospective clinical validation was performed, and corrected benchmark scores must not be read as estimates of clinical utility.

\section{Conclusion}
We asked whether public cervical-cytology classifiers are evaluated on genuinely independent cells, and found that released crop filenames conceal a consequential source-field hierarchy.
\cytofair{} recovers that hierarchy exactly, validates it by two independent routes, and measures what crop-level splitting shares.
On SIPaKMeD, \meanpm{92.62}{0.40}\% of test crops have a training parent and \meanpm{52.62}{1.97}\% share literal source pixels with a training crop; Cx22 and CRIC show the same construction-linked pattern; single-cell Herlev returns exactly zero and validates the instrument.
The shared signal is learnable at 109$\times$ chance and at 79$\times$ chance from inpainted background alone, locating it in acquisition context rather than cell morphology---which makes the exposure a property of the cropped-from-field construction itself.
Correcting the evaluation unit reduces accuracy for all six architectures, by a mean of 1.37 points, exceeding the 0.62-point spread between the two strongest published results, though five seeds cannot certify individual contrasts after correction.
The contribution is accordingly a benchmark correction and reusable measurement resource---exact provenance maps, fixed parent-disjoint splits, and the audit toolkit---rather than another incremental classifier.
Two extensions follow directly: a pooled omnibus endpoint with an adequately powered seed budget, replacing the per-architecture family recorded here as mis-specified; and the same grouping apparatus applied to provenance-aware instance segmentation on Cx22 and to two constructions whose grouping units are known but never enforced---specimen-disjoint multi-view Fish-Pak and temporally grouped Croatian Fish video---with deterministic manifests fixed before any model is trained.

\bibliographystyle{IEEEtran}
{\scriptsize\bibliography{references}}

\end{document}
